\documentclass[12pt,a4paper]{article}

\usepackage[utf8]{inputenc}
\usepackage[T1]{fontenc}
\usepackage[brazil]{babel}
\usepackage{geometry}
\usepackage{setspace}
\usepackage{hyperref}
\usepackage{titlesec}

\geometry{margin=2.5cm}
\onehalfspacing

\titleformat{\section}{\bfseries\large}{\thesection.}{0.5em}{}
\titleformat{\subsection}{\bfseries\normalsize}{\thesubsection.}{0.5em}{}

\title{\textbf{TERMO DE ABERTURA DO PROJETO (TAP)} \\ \vspace{0.5cm} \large Projeto PetSuam}
\author{Gabriel Ferreira Dantas \\ Igor da Silva Sant\textquotesingle anna \\ Ronivaldo Domingues de Andrade}
\date{Março de 2026}

\begin{document}

\maketitle

\section{Identificação do Projeto}
\textbf{Nome do Projeto:} PetSuam \\

\textbf{Curso:} Superior em Análise e Desenvolvimento de Sistemas \\

\textbf{Instituição:} Centro Universitário Augusto Motta -- UNISUAM \\

\textbf{Disciplina:} Projeto em Aplicativos Mobile \\

\textbf{Professor:} Bruno Santos Cezario \\

\textbf{Data de Início:} 19/03/2026 \\

\textbf{Previsão de Término:} 14/06/2026 \\

\section{Descrição do Projeto}
O projeto PetSuam consiste no desenvolvimento de um aplicativo mobile com o objetivo principal de funcionar como uma carteira de vacinação digital para animais de estimação.

A proposta visa centralizar informações importantes sobre a saúde dos pets, permitindo ao tutor acessar dados como histórico de vacinas, informações básicas do animal e facilitar o acompanhamento da saúde de forma prática e digital.

Inicialmente, o projeto contemplará as espécies caninas e felinas, com possibilidade de expansão futura para outras espécies em uma versão comercial.

\section{Objetivo do Projeto}
Desenvolver um aplicativo mobile funcional que permita:
\begin{itemize}
    \item Cadastro de tutores e seus pets;
    \item Armazenamento de informações básicas dos animais;
    \item Registro e acompanhamento de vacinas;
    \item Localização de pet shops próximos via integração com API de mapas;
\end{itemize}

\section{Justificativa}
Atualmente, muitos tutores não possuem controle organizado das vacinas de seus pets, utilizando métodos físicos ou desorganizados. O aplicativo busca resolver esse problema oferecendo uma solução digital acessível, prática e segura.

Além disso, o projeto apresenta potencial de expansão comercial, agregando funcionalidades que aumentam a segurança dos animais e facilitam a interação com clínicas veterinárias.

\section{Escopo do Projeto}

\subsection{Escopo do MVP}
\begin{itemize}
    \item Tela de login;
    \item Tela de cadastro do tutor;
    \item Cadastro de pets (cães e gatos);
    \item Listagem dos pets cadastrados;
    \item Tela de perfil do pet com:
    \begin{itemize}
        \item Dados básicos;
        \item Listagem de vacinas;
        \item Opção de adicionar novas vacinas;
    \end{itemize}
    \item Integração com API do Google Maps para localização de pet shops;
\end{itemize}

\subsection{Fora do Escopo (MVP)}
\begin{itemize}
    \item Cadastro de outras espécies;
    \item Sistema administrativo completo para clínicas;
    \item Funcionalidades avançadas de análise de dados;
\end{itemize}

\section{Funcionalidades Futuras (Versão Comercial)}
\begin{itemize}
    \item Geração de QR Code para cada pet;
    \item Página web pública acessível via QR Code com:
    \begin{itemize}
        \item Dados do tutor;
        \item Informações de segurança do animal;
    \end{itemize}
    \item Aplicativo administrativo para clínicas veterinárias;
    \item Registro automatizado de procedimentos e vacinas;
    \item Geração de dados estatísticos para estudos futuros;
\end{itemize}

\section{Público-Alvo}
O público-alvo do aplicativo são pessoas que possuem ao menos um animal de estimação, inicialmente focado em donos de cães e gatos.

\section{Tecnologias Utilizadas}
\begin{itemize}
    \item React Native (desenvolvimento mobile);
    \item API Google Maps;
    \item Integração futura com banco de dados remoto;
\end{itemize}

\section{Stakeholders}
\begin{itemize}
    \item Desenvolvedores: Gabriel Ferreira Dantas, Igor da Silva Sant\textquotesingle anna, Ronivaldo Domingues de Andrade;
    \item Instituição: UNISUAM;
    \item Professores da disciplina;
    \item Usuários finais (tutores de pets);
    \item Clínicas veterinárias (futuro);
\end{itemize}

\section{Cronograma Resumido}
\begin{itemize}
    \item Planejamento: Março de 2026;
    \item Desenvolvimento do MVP: Março a Maio de 2026;
    \item Testes: Maio de 2026;
    \item Entrega Final: Junho de 2026;
\end{itemize}

\section{Riscos do Projeto}
\begin{itemize}
    \item Dificuldade de aprendizado do React Native;
    \item Problemas na integração com APIs externas;
    \item Limitação de tempo para desenvolvimento completo;
\end{itemize}

\section{Links do Projeto}
\textbf{Protótipo (Figma):} \\
\url{https://www.figma.com/design/Y94pudhuCsMPn7SRlNW4ur/App-PetSuam?node-id=0-1&t=IjEc5wA5R9lsQVaO-1}

\vspace{0.3cm}

\noindent
\textbf{Repositório (GitHub):} \\
\url{https://github.com/adscolabhub/petsuam-app}

\section{Aprovação do Projeto}
Este documento formaliza a abertura do projeto PetSuam, autorizando o início de seu desenvolvimento conforme os termos aqui descritos.

\vspace{2cm}

\noindent
\begin{tabular}{ccc}
\rule{5cm}{0.4pt} & \rule{5cm}{0.4pt} & \rule{5cm}{0.4pt} \\
Gabriel Ferreira Dantas & Igor da Silva Sant\textquotesingle anna & Ronivaldo D. Andrade \\
\end{tabular}

\end{document}